\subsection{Problem S2-T3: explicit finite endpoint cells}

Use independent variables
\[
 (r,a,d,\tau,\beta)\in\mathbb R^5.
\]
Define
\[
 w=1-r,\qquad
 e=\sqrt{1-rw},\qquad
 \eta=1-e,\qquad
 p=e(1-e),\qquad
 k=\eta+a+d,
\]
\[
 s_R=\frac{k}{r},\qquad
 t_R=w+d,\qquad
 s_L=\frac{k}{w},\qquad
 t_L=r+a,
\]
and
\[
 z=\sqrt{1-\tau+\tau^2},\qquad
 b_-=\beta,\qquad
 b_+=z-\beta,
\]
\[
 q=1-z+\beta,\qquad
 c_\star=\frac q\tau,\qquad
 u_\star=1+\beta-\tau,
\]
\[
 \gamma_1=\frac d r,\qquad
 \gamma_5=\frac a w.
\]

\begin{definition}[Exact algebraic source domain]
\label{def:s2-t3-source-domain}
For \(\varepsilon\in\{1,2\}\), define
\[
 \mathcal B_{\rm T3}^{\varepsilon}
 =
 \left\{
 (r,a,d,\tau,\beta):
 \begin{array}{l}
 (r,a,d)\in\mathcal D_{\rm C\varepsilon},\\
 0<\tau<1,\quad
 0<\beta<\dfrac{\tau z}{1+z},\\
 0<b_-<1,\quad 0<b_+<1,\\
 0<q<\dfrac12,\quad
 0<c_\star<u_\star<1
 \end{array}
 \right\},
\]
where
\[
 \mathcal D_{\rm C\varepsilon}
 =
 \begin{cases}
 \mathcal D_{\rm C1},&\varepsilon=1,\\
 \mathcal D_{\rm C2},&\varepsilon=2.
 \end{cases}
\]
The variables \((\tau,\beta)\) are an exact radical parameterization; no
translated-set or containment predicate occurs in this domain.
\end{definition}

Introduce endpoint variables
\[
 A_1,A_5,C_1,C_5\in[0,1].
\]
The right-end residual cells are
\[
 \begin{aligned}
 \mathcal R_0&:\quad b_+<s_R,\quad A_1=q,\\
 \mathcal R_1&:\quad s_R\le b_+<t_R,\quad A_1=r-d,\\
 \mathcal R_2&:\quad t_R\le b_+,\quad A_1=q.
 \end{aligned}
\]
The left-end cells are
\[
 \mathcal L_{1,0}:\quad A_5=1-\beta
\]
on \(\mathcal D_{\rm C1}\), and on \(\mathcal D_{\rm C2}\),
\[
 \begin{aligned}
 \mathcal L_{2,0}&:\quad \beta<s_L,\quad A_5=1-\beta,\\
 \mathcal L_{2,1}&:\quad s_L\le\beta<t_L,\quad A_5=1-t_L,\\
 \mathcal L_{2,2}&:\quad t_L\le\beta,\quad A_5=1-\beta.
 \end{aligned}
\]
The radial cells are
\[
 \begin{aligned}
 \mathcal K_0&:\quad
 1-u_\star\le\gamma_1\le1-c_\star,\quad C_1=c_\star,\\
 \mathcal K_1&:\quad
 \gamma_1<1-u_\star,\quad C_1=1-\gamma_1,\\
 \mathcal K_2&:\quad
 1-c_\star<\gamma_1,\quad C_1=1-\gamma_1.
 \end{aligned}
\]
In every cell,
\[
 C_5=1-\gamma_5.
\]
Equality at either endpoint of the hit interval belongs to \(\mathcal K_0\).

For \(\varepsilon=1\), put \(J_\varepsilon=\{0\}\); for
\(\varepsilon=2\), put \(J_\varepsilon=\{0,1,2\}\).

\begin{definition}[Finite T3 cell union]
\label{def:s2-t3-finite-union}
For
\[
 \varepsilon\in\{1,2\},\quad
 i,h\in\{0,1,2\},\quad
 j\in J_\varepsilon,
\]
let \(\Omega_{\rm T3}^{\varepsilon,i,j,h}\) be the set of tuples
\[
 \xi=(r,a,d,\tau,\beta,A_1,A_5,C_1,C_5)
\]
satisfying
\[
 \mathcal B_{\rm T3}^{\varepsilon},
 \qquad
 \mathcal R_i,\qquad
 \mathcal L_{\varepsilon,j},\qquad
 \mathcal K_h,\qquad
 C_5=1-\gamma_5,
\]
and
\[
 \frac12<C_1<1,\qquad
 \frac12<C_5<1.
\]
Define the easy subcell by adding
\[
 A_1+A_5>1.
\]
For \((\ell_5,\ell_1)\in\mathfrak L^2\), define the hard subcell by adding
\[
 A_1+A_5\le1,\qquad
 \mathrm{Cell}_{\ell_5}(C_5,A_5),\qquad
 \mathrm{Cell}_{\ell_1}(C_1,A_1).
\]
Finally,
\[
 \begin{aligned}
 \Omega_{\rm T3}
 ={}&
 \bigcup_{\substack{\varepsilon,i,j,h}}
 \Omega_{\rm T3,easy}^{\varepsilon,i,j,h}\\
 &{}\cup
 \bigcup_{\substack{\varepsilon,i,j,h\\
                     (\ell_5,\ell_1)\in\mathfrak L^2}}
 \Omega_{\rm T3,hard}^{\varepsilon,i,j,h,\ell_5,\ell_1}.
 \end{aligned}
\]
Both unions are finite.  Every cell is specified by explicit real equalities,
strict or weak inequalities, and the displayed radicals.
\end{definition}

\begin{definition}[Problem S2-T3]
\label{prob:s2-t3-endpoint}
For
\[
 \xi=(r,a,d,\tau,\beta,A_1,A_5,C_1,C_5)\in\Omega_{\rm T3},
\]
define
\[
 \Psi_{\rm T3}(\xi)
 =F(C_5,A_5)+F(C_1,A_1)-1.
\]
The optimization problem is
\[
 \boxed{
 \Psi_{\rm T3}(\xi)<0
 \qquad
 (\xi\in\Omega_{\rm T3}).
 }
\]
\end{definition}

On the easy cells the sign follows from \(F(c,x)\le1-x\).  The hard cells
are exactly the finite four-label inventory.  No contact selector, equality
assignment, radical sign, residual order, or radial hit/miss alternative is
left implicit.
